\documentclass[11pt]{article}

\usepackage[utf8]{inputenc}
\usepackage[T1]{fontenc}
\usepackage{lmodern}
\usepackage{geometry}
\usepackage{amsmath,amssymb,amsfonts,amsthm,mathtools}
\usepackage{bm}
\usepackage{enumitem}
\usepackage{microtype}
\usepackage{hyperref}
\usepackage{titlesec}
\usepackage{booktabs}
\usepackage{array}
\usepackage{setspace}
\usepackage{mathrsfs}
\usepackage{physics}

\geometry{margin=1in}
\hypersetup{colorlinks=true, linkcolor=black, urlcolor=blue, citecolor=black}
\setlist[enumerate]{topsep=0.4em,itemsep=0.35em}
\setlist[itemize]{topsep=0.3em,itemsep=0.25em}
\titleformat{\section}{\large\bfseries}{\thesection.}{0.5em}{}
\titleformat{\subsection}{\normalsize\bfseries}{\thesubsection.}{0.5em}{}
\setstretch{1.06}

\newcommand{\Aether}{\AE{}ther}
\newcommand{\Aflow}{\AE{}ther-flow}
\newcommand{\TheoryName}{\Aflow{} Interpretation of Relativity}
\newcommand{\TheoryProgram}{\Aether{} / \Aflow{} framework}

\newcommand{\CompletedMetric}{g^{\sharp}}
\newcommand{\MatterSpace}{\Psi}

\newcommand{\Car}{\mathrm{car}}
\newcommand{\Cur}{\mathrm{cur}}
\newcommand{\Hom}{\mathrm{hom}}

\newcommand{\ForSpace}[1]{\mathcal I_{#1}^{\mathrm{for}}}
\newcommand{\PrimShell}{\mathcal S_{\mathrm{prim}}}
\newcommand{\MetricOut}{\mathscr G_{\mathrm{out}}}
\newcommand{\MatterOut}{\mathscr T_{\mathrm{out}}}

\newcommand{\ProtoSectionBody}[1]{\mathcal B_{#1}^{\mathrm{psec}}}
\newcommand{\ProtoSectionShell}[1]{\mathcal Z_{#1}^{\mathrm{psec}}}

\newcommand{\PrePreProtoSectionPullbackCore}{\mathfrak S^{\mathrm{sub,proto,pppresec,pb}}}
\newcommand{\PrePreProtoSectionVertical}[1]{V_{#1}^{\mathrm{pppresec}}}
\newcommand{\PrePreProtoSectionResidualBody}[1]{\mathcal D_{#1}^{\mathrm{pppresec}}}

\newcommand{\PullbackOperatorCore}{\mathfrak S^{\mathrm{sub,proto,pppresec,pb,op}}}
\newcommand{\PullbackBodyOperatorCore}{\mathfrak S^{\mathrm{sub,proto,pppresec,pb,bodyop}}}
\newcommand{\PullbackBodyResponseCore}{\mathfrak S^{\mathrm{sub,proto,pppresec,pb,bodyresp}}}
\newcommand{\PreProtoSectionResidualPackage}{\mathfrak R^{\mathrm{sub,proto,presec,res}}}
\newcommand{\PreProtoSectionPackage}{\mathfrak R^{\mathrm{sub,proto,presec}}}
\newcommand{\PreProtoSectionMinSectionCore}{\mathfrak R^{\mathrm{sub,proto,presec,sec}}}

\newcommand{\SubstrateCarrier}[1]{U_{#1}^{\mathrm{disp,resp}}}
\newcommand{\SubstrateProjection}[1]{P_{#1}^{\mathrm{disp,resp}}}
\newcommand{\SubstrateOperator}[1]{K_{#1}^{\mathrm{disp,resp}}}
\newcommand{\SubstrateKernel}[1]{N_{#1}^{\mathrm{disp,resp}}}
\newcommand{\SchurResponse}[1]{M_{#1}^{\mathrm{disp,resp}}}
\newcommand{\CanonicalLift}[1]{J_{#1}^{\mathrm{disp,resp,can}}}
\newcommand{\CoreForcedBodyResponse}[1]{\widehat{\mathscr R}_{#1}^{\mathrm{disp,resp,projop}}}
\newcommand{\CanonicalImageBody}[1]{\mathcal U_{#1}^{\mathrm{disp,resp,cbimg}}}
\newcommand{\CanonicalImageShell}[1]{\mathcal Z_{#1}^{\mathrm{disp,resp,cbimg}}}

\newcommand{\CompatibleImageBody}[1]{\widetilde{\mathcal U}_{#1}^{\mathrm{disp,resp,cbimg}}}
\newcommand{\CompatibleImageLift}[1]{\widetilde J_{#1}^{\mathrm{disp,resp,cbimg}}}
\newcommand{\CompatibleImageResponse}[1]{\widetilde{\mathscr R}_{#1}^{\mathrm{disp,resp,cbimg}}}

\newtheorem{definition}{Definition}
\newtheorem{proposition}{Proposition}
\newtheorem{remark}{Remark}
\newtheorem{corollary}{Corollary}
\newtheorem{theorem}{Theorem}

\title{\textbf{The \TheoryName{}}\\[0.4em]
\large Primitive-Reservoir Observer-Reduction\\
\large Section-Centered Hidden-Response\\
\large Canonical Image-Body Necessity / Uniqueness Theorem}
\author{}
\date{}

\begin{document}

\maketitle

\begin{abstract}
The current active-primary theorem already compressed the full section-centered hidden-response projection-operator-core frontier to the smaller canonical current-body image core by showing that off-image kernel directions are benchmark neutral once the current displacement body and the coercive response operator are fixed. The routing note below that frontier then made the next honest options explicit: either derive the smaller image core from still more primitive explicit substrate structure, or, if that full derivation is not yet clean, prove a sharper forcing or uniqueness theorem for a nontrivial constituent or relation inside that smaller core.

The present note takes that honest internal route. It proves that the canonical image body itself is already uniquely forced at benchmark scope once the fixed current displacement body, the fixed surjective projection, the fixed coercive positive self-adjoint substrate-response operator, the fixed exact-shell image, and the fixed current body-response route are held fixed. More precisely, any benchmark-faithful body candidate inside the common carrier whose restricted projection bijects onto the fixed current displacement body and whose induced quadratic response agrees with the already forced current response must coincide fiberwise with the canonical minimal-lift image. There is therefore no remaining benchmark-facing body-level freedom inside the current-body-image-core frontier.

The benchmark boundary remains fixed throughout. The one operative metric is still exactly \(\CompletedMetric\), the ordinary matter route is unchanged, there is no second operative metric, the low-energy matter law is not enlarged, and no stronger promotion language is licensed. The positive result is therefore a genuine sharper theorem strictly inside the current image-core frontier: the body constituent is no longer independently live there. What remains open is to derive or uniquely force the restricted projection / canonical-lift relation, the exact-shell image, the response operator restricted to the canonical image, or the full image core itself from still more primitive explicit substrate structure.
\end{abstract}

\section{Introduction}

The active derivational program of \TheoryName{} again reaches a stage where depth alone is no longer the honest default move. The preceding current theorem already proved that the full projection-operator-core body presentation is not the minimal benchmark-facing datum on the section-centered hidden-response route. Once the fixed current displacement body, the surjective projection, and the coercive response operator are held fixed, every point of the larger body splits uniquely into the canonical minimal-lift representative of the same current displacement plus a benchmark-neutral kernel direction. The live burden therefore moved from the larger body presentation to the smaller canonical current-body image core.

That still leaves an internal question. Inside the current image core itself, does the image body remain an independently live benchmark-facing constituent, or is it already forced once the remaining route data are fixed? The present note proves the stronger second statement. On the recorded route, any body candidate that projects bijectively to the fixed current displacement body and induces the already forced current quadratic response must coincide with the canonical minimal-lift image. The body constituent therefore carries no remaining benchmark-facing presentational freedom at current scope.

\paragraph{Framework and claim boundary.}
\TheoryProgram{} denotes the broader \Aether{} / \Aflow{} framework built on the claims that the \Aether{} is the underlying four-dimensional substrate of reality, the \Aflow{} is its intrinsic ordered motion, observed three-dimensional space is the local experiential slice of that deeper substrate, S-time is the experienced order of change arising from matter, light, and the \Aflow{}, and observed expansion is the three-dimensional appearance of deeper four-dimensional ordered motion. Within that framework, \TheoryName{} denotes the exact-closure theory in which the effective gravitational dynamics are adopted to be exactly Einsteinian with universal matter coupling. In the active sequence, \emph{adoption} means use of that established relativistic dynamics without claiming substrate derivation, while \emph{derivation} is reserved for a first-principles recovery from explicit substrate variables.

The present note stays strictly inside that boundary. It does not introduce a second operative metric, it does not enlarge the low-energy matter law, and it does not reopen the already-screened full projection-operator-core body presentation, the full displacement-response substrate law, or the larger pullback-body-operator, pullback-operator, pullback-quadratic, hidden-response, residual-normal-form, or pre-proto-section presentations as if they were again the live burden. Its narrower benchmark-facing role is to prove that the current image-body constituent itself is already uniquely forced on the recorded route.

\section{Fixed Inputs}

\begin{definition}[Recorded primitive continuation data]
\label{def:fixed}
Fix the following continuation data from the active primitive-reservoir observer-reduction line.
\begin{enumerate}
    \item For each sector label \(\sigma\in\{\Car,\Cur,\Hom\}\), a primitive forcing shell
    \[
    \ForSpace{\sigma}.
    \]
    \item The product shell
    \[
    \PrimShell
    \equiv
    \ForSpace{\Car}\times \ForSpace{\Cur}\times \ForSpace{\Hom}.
    \]
    \item The recorded metric and ordinary-matter outputs
    \[
    \MetricOut:\PrimShell\to\{\text{observer metrics}\},
    \qquad
    \MatterOut:\PrimShell\times\MatterSpace\to\{\text{observer stress tensors}\},
    \]
    satisfying
    \begin{equation}
    \MetricOut(\mathbf w)=\CompletedMetric,
    \qquad
    \MatterOut(\mathbf w,\psi)
    =
    T_{\mu\nu}^{\mathrm{matter}}[\CompletedMetric,\psi]
    \label{eq:fixedoutputs}
    \end{equation}
    for every \(\mathbf w\in\PrimShell\) and every matter configuration \(\psi\in\MatterSpace\).
\end{enumerate}
\end{definition}

\begin{definition}[Recorded proto-section benchmark base and current image-body forcing data]
\label{def:route}
Fix \(\sigma\in\{\Car,\Cur,\Hom\}\). The active line already fixes:
\begin{enumerate}
    \item the current proto-section benchmark base \(\ProtoSectionBody{\sigma}\) and exact proto-section shell \(\ProtoSectionShell{\sigma}\);
    \item the current section-centered displacement body
    \[
    \PrePreProtoSectionResidualBody{\sigma}
    \subset
    \ProtoSectionBody{\sigma}\times\PrePreProtoSectionVertical{\sigma}
    \]
    whose fiber
    \[
    \PrePreProtoSectionResidualBody{\sigma}(y)
    \equiv
    \left\{
    v\in\PrePreProtoSectionVertical{\sigma}:(y,v)\in\PrePreProtoSectionResidualBody{\sigma}
    \right\}
    \]
    is nonempty, compact, convex, contains \(0\), and has nonempty interior for every \(y\in\ProtoSectionBody{\sigma}\);
    \item a finite-dimensional Euclidean substrate carrier
    \[
    \SubstrateCarrier{\sigma};
    \]
    \item a smooth family of surjective linear displacement projections
    \[
    \SubstrateProjection{\sigma}(y):
    \SubstrateCarrier{\sigma}\to\PrePreProtoSectionVertical{\sigma},
    \qquad
    y\in\ProtoSectionBody{\sigma};
    \]
    \item a smooth family of positive self-adjoint substrate-response operators
    \[
    \SubstrateOperator{\sigma}(y):
    \SubstrateCarrier{\sigma}\to\SubstrateCarrier{\sigma},
    \qquad
    y\in\ProtoSectionBody{\sigma},
    \]
    satisfying the uniform coercive bound
    \begin{equation}
    \ev{u,\SubstrateOperator{\sigma}(y)u}
    \ge
    \mu_{\sigma}\,\|u\|^2
    \qquad
    \text{for all }
    y\in\ProtoSectionBody{\sigma},
    \ \ u\in\SubstrateCarrier{\sigma}
    \label{eq:coercive}
    \end{equation}
    for some constant \(\mu_{\sigma}>0\);
    \item the fixed exact-shell/current-body compatibility relation
    \[
    \left\{
    y\in\ProtoSectionBody{\sigma}:0\in\PrePreProtoSectionResidualBody{\sigma}(y)
    \right\}
    =
    \ProtoSectionShell{\sigma}
    \]
    and the fixed exact-shell image
    \begin{equation}
    \CanonicalImageShell{\sigma}
    \equiv
    \left\{
    (y,0):y\in\ProtoSectionShell{\sigma}
    \right\};
    \label{eq:imageshell}
    \end{equation}
    \item benchmark faithfulness, meaning preservation of \eqref{eq:fixedoutputs}, no second operative metric, no enlarged low-energy matter law, no change to the ordinary matter route, and no premature promotion from adoption to completed derivation.
\end{enumerate}
\end{definition}

\begin{definition}[Recorded downstream chain]
\label{def:downstream}
Fix \(\sigma\in\{\Car,\Cur,\Hom\}\). The active line already records that, once the same benchmark-faithful current body-response core on the fixed displacement body is available, the current observer-relevant pullback body-response core \(\PullbackBodyResponseCore\), the current pullback body-operator core \(\PullbackBodyOperatorCore\), the current pullback-operator core \(\PullbackOperatorCore\), the current pullback-quadratic core \(\PrePreProtoSectionPullbackCore\), the current section-centered residual-normal-form realization \(\PreProtoSectionResidualPackage\), the current benchmark-faithful pre-proto-section package \(\PreProtoSectionPackage\), the current minimizer-section core \(\PreProtoSectionMinSectionCore\), and the previously recorded continuity-Hessian compressed core and benchmark one-metric observer package follow unchanged.
\end{definition}

\section{Canonical Lift and Canonical Image Body}

\begin{proposition}[Positive Schur-response operator]
\label{prop:schur}
For each \(y\in\ProtoSectionBody{\sigma}\), define
\[
\SchurResponse{\sigma}(y)
\equiv
\SubstrateProjection{\sigma}(y)
\SubstrateOperator{\sigma}(y)^{-1}
\SubstrateProjection{\sigma}(y)^*:
\PrePreProtoSectionVertical{\sigma}\to\PrePreProtoSectionVertical{\sigma},
\]
where \((\cdot)^*\) denotes Euclidean adjoint with respect to the fixed inner products on \(\SubstrateCarrier{\sigma}\) and \(\PrePreProtoSectionVertical{\sigma}\). Then \(\SchurResponse{\sigma}(y)\) is positive self-adjoint and invertible for every \(y\in\ProtoSectionBody{\sigma}\).
\end{proposition}

\begin{proof}
Self-adjointness is immediate from the self-adjointness of \(\SubstrateOperator{\sigma}(y)^{-1}\). For any \(0\neq w\in\PrePreProtoSectionVertical{\sigma}\),
\[
\ev{w,\SchurResponse{\sigma}(y)w}
=
\ev{
\SubstrateOperator{\sigma}(y)^{-1/2}\SubstrateProjection{\sigma}(y)^*w,
\SubstrateOperator{\sigma}(y)^{-1/2}\SubstrateProjection{\sigma}(y)^*w
}
>0
\]
because \(\SubstrateProjection{\sigma}(y)\) is surjective and therefore \(\SubstrateProjection{\sigma}(y)^*\) is injective. Hence \(\SchurResponse{\sigma}(y)\) is positive definite and invertible.
\end{proof}

\begin{definition}[Canonical lift and core-forced response]
\label{def:canonical}
For each \(y\in\ProtoSectionBody{\sigma}\), define the \emph{canonical minimal lift}
\begin{equation}
\CanonicalLift{\sigma}(y)
\equiv
\SubstrateOperator{\sigma}(y)^{-1}\SubstrateProjection{\sigma}(y)^*
\SchurResponse{\sigma}(y)^{-1}
:
\PrePreProtoSectionVertical{\sigma}\to\SubstrateCarrier{\sigma}.
\label{eq:canonical}
\end{equation}
Define the quadratic substrate-response energy by
\[
\mathscr E_{\sigma}(y,u)
\equiv
\frac{1}{2}
\ev{u,\SubstrateOperator{\sigma}(y)u}
\]
and the \emph{core-forced body response}
\begin{equation}
\CoreForcedBodyResponse{\sigma}(y,v)
\equiv
\frac{1}{2}
\ev{
\CanonicalLift{\sigma}(y)v,
\SubstrateOperator{\sigma}(y)\CanonicalLift{\sigma}(y)v
}
\qquad
\text{for }
(y,v)\in\PrePreProtoSectionResidualBody{\sigma}.
\label{eq:coreresponse}
\end{equation}
\end{definition}

\begin{proposition}[Canonical lift properties]
\label{prop:canonical}
For every \(y\in\ProtoSectionBody{\sigma}\),
\begin{enumerate}
    \item
    \[
    \SubstrateProjection{\sigma}(y)\circ\CanonicalLift{\sigma}(y)
    =
    \mathrm{id}_{\PrePreProtoSectionVertical{\sigma}};
    \]
    \item
    \[
    \ev{
    \CanonicalLift{\sigma}(y)v,
    \SubstrateOperator{\sigma}(y)k
    }
    =
    0
    \]
    for every \(v\in\PrePreProtoSectionVertical{\sigma}\) and every \(k\in\SubstrateKernel{\sigma}(y)\equiv\ker\SubstrateProjection{\sigma}(y)\);
    \item every \(u\in\SubstrateCarrier{\sigma}\) admits the unique decomposition
    \[
    u
    =
    \CanonicalLift{\sigma}(y)\bigl(\SubstrateProjection{\sigma}(y)u\bigr)
    +
    k,
    \qquad
    k\in\SubstrateKernel{\sigma}(y);
    \]
    \item \(\CanonicalLift{\sigma}(y)\) is linear in the displacement argument and smooth in \(y\).
\end{enumerate}
\end{proposition}

\begin{proof}
Item (1) follows directly from Definition \ref{def:canonical}. For item (2), let \(k\in\SubstrateKernel{\sigma}(y)\). Then
\[
\ev{
\CanonicalLift{\sigma}(y)v,
\SubstrateOperator{\sigma}(y)k
}
=
\ev{
\SubstrateProjection{\sigma}(y)^*\SchurResponse{\sigma}(y)^{-1}v,
k
}
=
\ev{
\SchurResponse{\sigma}(y)^{-1}v,
\SubstrateProjection{\sigma}(y)k
}
=
0.
\]
Item (3) follows by setting
\[
k\equiv u-\CanonicalLift{\sigma}(y)\bigl(\SubstrateProjection{\sigma}(y)u\bigr)
\]
and applying item (1). Uniqueness follows because item (1) fixes the non-kernel part. Linearity and smoothness are immediate from \eqref{eq:canonical}.
\end{proof}

\begin{proposition}[Unique energy-minimizing lift]
\label{prop:min}
For every \(y\in\ProtoSectionBody{\sigma}\) and every \(v\in\PrePreProtoSectionVertical{\sigma}\), the canonical lift \(\CanonicalLift{\sigma}(y)v\) is the unique minimizer of
\begin{equation}
\min\left\{
\mathscr E_{\sigma}(y,u):
u\in\SubstrateCarrier{\sigma},
\ \SubstrateProjection{\sigma}(y)u=v
\right\}.
\label{eq:minproblem}
\end{equation}
\end{proposition}

\begin{proof}
Let \(u\in\SubstrateCarrier{\sigma}\) satisfy \(\SubstrateProjection{\sigma}(y)u=v\). By Proposition \ref{prop:canonical}(3),
\[
u=\CanonicalLift{\sigma}(y)v+k,
\qquad
k\in\SubstrateKernel{\sigma}(y).
\]
Using Proposition \ref{prop:canonical}(2),
\[
\ev{u,\SubstrateOperator{\sigma}(y)u}
=
\ev{
\CanonicalLift{\sigma}(y)v,
\SubstrateOperator{\sigma}(y)\CanonicalLift{\sigma}(y)v
}
+
\ev{k,\SubstrateOperator{\sigma}(y)k}.
\]
The last term is nonnegative by positivity of \(\SubstrateOperator{\sigma}(y)\), and by the coercive bound \eqref{eq:coercive} it vanishes only if \(k=0\). Therefore the unique minimizer of \eqref{eq:minproblem} is \(\CanonicalLift{\sigma}(y)v\).
\end{proof}

\begin{definition}[Canonical image body]
\label{def:image}
For each \(y\in\ProtoSectionBody{\sigma}\), define the \emph{canonical image body}
\begin{equation}
\CanonicalImageBody{\sigma}(y)
\equiv
\CanonicalLift{\sigma}(y)\bigl(\PrePreProtoSectionResidualBody{\sigma}(y)\bigr)
\subset
\SubstrateCarrier{\sigma}.
\label{eq:imagebody}
\end{equation}
\end{definition}

\begin{remark}
Definition \ref{def:image} is the current body-level constituent singled out by the preceding image-core collapse theorem. The present note asks whether even this body constituent still carries independent benchmark-facing freedom once the remaining route data are held fixed.
\end{remark}

\section{Compatible Benchmark-Facing Body Candidates}

\begin{definition}[Benchmark-faithful compatible current-image body]
\label{def:compatible}
A \emph{benchmark-faithful compatible current-image body} for sector \(\sigma\) consists of fiber sets
\[
\CompatibleImageBody{\sigma}(y)\subset\SubstrateCarrier{\sigma},
\qquad
y\in\ProtoSectionBody{\sigma},
\]
together with inverse sections
\[
\CompatibleImageLift{\sigma}(y):
\PrePreProtoSectionResidualBody{\sigma}(y)\to\CompatibleImageBody{\sigma}(y),
\]
such that:
\begin{enumerate}
    \item each fiber \(\CompatibleImageBody{\sigma}(y)\) is nonempty, compact, and convex;
    \item the restricted projection
    \[
    \SubstrateProjection{\sigma}(y)\big|_{\CompatibleImageBody{\sigma}(y)}
    :
    \CompatibleImageBody{\sigma}(y)\to\PrePreProtoSectionResidualBody{\sigma}(y)
    \]
    is a bijection with inverse \(\CompatibleImageLift{\sigma}(y)\);
    \item the exact-shell image is fixed:
    \[
    \left\{
    \bigl(y,\CompatibleImageLift{\sigma}(y)0\bigr):
    y\in\ProtoSectionShell{\sigma}
    \right\}
    =
    \CanonicalImageShell{\sigma};
    \]
    \item the induced response
    \begin{equation}
    \CompatibleImageResponse{\sigma}(y,v)
    \equiv
    \frac{1}{2}
    \ev{
    \CompatibleImageLift{\sigma}(y)v,
    \SubstrateOperator{\sigma}(y)\CompatibleImageLift{\sigma}(y)v
    }
    \label{eq:compatresponse}
    \end{equation}
    agrees with the fixed current body-response route:
    \[
    \CompatibleImageResponse{\sigma}(y,v)
    =
    \CoreForcedBodyResponse{\sigma}(y,v)
    \qquad
    \text{for every }
    (y,v)\in\PrePreProtoSectionResidualBody{\sigma};
    \]
    \item benchmark faithfulness in the sense of Definition \ref{def:route}(7).
\end{enumerate}
\end{definition}

\begin{remark}
Definition \ref{def:compatible} is intentionally internal to the current image-core frontier. It does not yet claim a deeper explicit substrate law below the image core. It asks only whether a different benchmark-facing body constituent can remain once the fixed displacement body and the fixed projection / response data are held fixed.
\end{remark}

\begin{proposition}[The canonical image body is compatible]
\label{prop:exists}
The canonical image body of Definition \ref{def:image} is a benchmark-faithful compatible current-image body with inverse section \(\CanonicalLift{\sigma}(y)\big|_{\PrePreProtoSectionResidualBody{\sigma}(y)}\).
\end{proposition}

\begin{proof}
Each fiber \(\CanonicalImageBody{\sigma}(y)\) is nonempty, compact, and convex because it is the linear image of the nonempty compact convex fiber \(\PrePreProtoSectionResidualBody{\sigma}(y)\). Proposition \ref{prop:canonical}(1) shows that the restricted projection is inverse to the canonical lift on the fixed displacement body. By linearity, \(\CanonicalLift{\sigma}(y)0=0\), so item (3) of Definition \ref{def:compatible} is exactly \eqref{eq:imageshell}. Item (4) is Definition \ref{def:canonical}. Item (5) holds because the fixed displacement body together with the same body-response route preserve the same benchmark-facing downstream package by Definitions \ref{def:route}(7) and \ref{def:downstream}.
\end{proof}

\begin{proposition}[Compatible lifts are canonical]
\label{prop:lift}
Let \(\CompatibleImageBody{\sigma}\) be any benchmark-faithful compatible current-image body. Then
\[
\CompatibleImageLift{\sigma}(y)v
=
\CanonicalLift{\sigma}(y)v
\]
for every \(y\in\ProtoSectionBody{\sigma}\) and every \(v\in\PrePreProtoSectionResidualBody{\sigma}(y)\).
\end{proposition}

\begin{proof}
Fix \(y\in\ProtoSectionBody{\sigma}\) and \(v\in\PrePreProtoSectionResidualBody{\sigma}(y)\). By Definition \ref{def:compatible}(2), the point \(\CompatibleImageLift{\sigma}(y)v\) satisfies
\[
\SubstrateProjection{\sigma}(y)\CompatibleImageLift{\sigma}(y)v=v.
\]
Hence it is an admissible point in the constrained minimization problem \eqref{eq:minproblem}. By Definition \ref{def:compatible}(4),
\[
\frac{1}{2}
\ev{
\CompatibleImageLift{\sigma}(y)v,
\SubstrateOperator{\sigma}(y)\CompatibleImageLift{\sigma}(y)v
}
=
\CoreForcedBodyResponse{\sigma}(y,v)
=
\frac{1}{2}
\ev{
\CanonicalLift{\sigma}(y)v,
\SubstrateOperator{\sigma}(y)\CanonicalLift{\sigma}(y)v
}.
\]
Therefore \(\CompatibleImageLift{\sigma}(y)v\) achieves the same minimal value as the canonical lift. Proposition \ref{prop:min} says that minimizer is unique. Hence \(\CompatibleImageLift{\sigma}(y)v=\CanonicalLift{\sigma}(y)v\).
\end{proof}

\begin{corollary}[Compatible bodies are forced]
\label{cor:body}
If \(\CompatibleImageBody{\sigma}\) is a benchmark-faithful compatible current-image body, then for every \(y\in\ProtoSectionBody{\sigma}\),
\[
\CompatibleImageBody{\sigma}(y)
=
\CanonicalImageBody{\sigma}(y).
\]
Moreover, the restricted projection on that body is inverse to the canonical lift and the exact-shell image is exactly \(\CanonicalImageShell{\sigma}\).
\end{corollary}

\begin{proof}
By Definition \ref{def:compatible}(2),
\[
\CompatibleImageBody{\sigma}(y)
=
\CompatibleImageLift{\sigma}(y)\bigl(\PrePreProtoSectionResidualBody{\sigma}(y)\bigr).
\]
Proposition \ref{prop:lift} identifies this lift with \(\CanonicalLift{\sigma}(y)\), and Definition \ref{def:image} then gives the stated body equality. The remaining claims follow immediately from Definition \ref{def:compatible}(2)--(3) and Proposition \ref{prop:canonical}(1).
\end{proof}

\begin{proposition}[The canonical image body is sufficient at current benchmark scope]
\label{prop:sufficient}
The canonical image body together with the fixed current body-response route is sufficient to determine every benchmark-facing downstream statement currently available below the image-core frontier.
\end{proposition}

\begin{proof}
By Proposition \ref{prop:exists}, the canonical image body realizes the same current body-response route \(\CoreForcedBodyResponse{\sigma}\) on the fixed displacement body. Definition \ref{def:downstream} then gives the same current pullback body-response core \(\PullbackBodyResponseCore\), the same pullback body-operator core \(\PullbackBodyOperatorCore\), the same pullback-operator core \(\PullbackOperatorCore\), the same pullback-quadratic core \(\PrePreProtoSectionPullbackCore\), the same residual-normal-form realization \(\PreProtoSectionResidualPackage\), the same benchmark-faithful pre-proto-section package \(\PreProtoSectionPackage\), the same minimizer-section core \(\PreProtoSectionMinSectionCore\), and the same downstream benchmark one-metric observer package.
\end{proof}

\section{Main Theorem}

\begin{theorem}[Canonical image-body necessity / uniqueness on the current route]
\label{thm:main}
Fix the primitive continuation data of Definition \ref{def:fixed}, the recorded proto-section benchmark base and current image-body forcing data of Definition \ref{def:route}, and the recorded downstream chain of Definition \ref{def:downstream}. Then the canonical image body \(\CanonicalImageBody{\sigma}\) is the unique benchmark-facing body compatible with the fixed current displacement body, the fixed surjective projection, the fixed coercive response operator, the fixed exact-shell image, and the fixed current body-response route. More precisely:
\begin{enumerate}
    \item every benchmark-faithful compatible current-image body has inverse section equal to the canonical lift by Proposition \ref{prop:lift};
    \item consequently every such compatible body equals the canonical image body fiberwise and carries the same restricted projection / canonical-lift relation and the same exact-shell image by Corollary \ref{cor:body};
    \item the canonical image body is sufficient to determine every currently available benchmark-facing downstream consequence by Proposition \ref{prop:sufficient};
    \item therefore no genuinely different benchmark-facing body constituent remains live inside the current current-body-image-core frontier.
\end{enumerate}
\end{theorem}

\begin{proof}
Item (1) is Proposition \ref{prop:lift}. Item (2) is Corollary \ref{cor:body}. Item (3) is Proposition \ref{prop:sufficient}. Item (4) is the routing consequence of items (1)--(3): any alternative compatible body is not genuinely alternative because it equals \(\CanonicalImageBody{\sigma}\), while any body that changes the fixed response route, the fixed shell image, or the downstream benchmark package is not benchmark-faithfully compatible with the current route.
\end{proof}

\begin{corollary}[No genuinely different body datum remains at current scope]
\label{cor:minimal}
Inside the present section-centered hidden-response current-body-image-core frontier, there is no genuinely different benchmark-facing body datum below \(\CanonicalImageBody{\sigma}\) that is compatible with the fixed current route data.
\end{corollary}

\begin{proof}
By Theorem \ref{thm:main}, every benchmark-faithful compatible body coincides with \(\CanonicalImageBody{\sigma}\). Hence no distinct compatible body datum remains live at current benchmark scope.
\end{proof}

\begin{corollary}[Primary post-uniqueness burden]
\label{cor:next}
After Theorem \ref{thm:main}, the next honest benchmark-facing manuscript below the present frontier must do at least one of the following:
\begin{enumerate}
    \item derive or force the restricted projection / canonical-lift relation on the fixed current displacement body from still more primitive explicit substrate structure;
    \item derive or force the exact-shell image or the response operator restricted to the canonical image from still more primitive explicit substrate structure;
    \item derive or force the full canonical current-body image core in one cleaner deeper package if that unified derivation is genuinely available;
    \item prove a sharper compatibility, necessity, or uniqueness theorem for some remaining nontrivial constituent or relation inside the current image-core data;
    \item do one of the previous four while preserving the same benchmark one-metric package, the same ordinary matter route, and the same claim boundary.
\end{enumerate}
Another same-output deeper relay that merely preserves the same canonical image body, another restoration of the full substrate-body presentation, another replay of the full projection-operator core, or another replay of the full displacement-response substrate law is not the honest default next move.
\end{corollary}

\begin{proof}
Theorem \ref{thm:main} and Corollary \ref{cor:minimal} show that the body constituent itself is already spent at benchmark scope on the recorded route. Therefore a new primary manuscript must either derive or sharpen some remaining nontrivial constituent or relation inside the current image-core data, or derive the whole image core itself from still more primitive explicit substrate structure. If it does not, then it preserves the same already-forced body datum and is benchmark neutral at project-status level.
\end{proof}

\begin{proposition}[The benchmark package remains fixed]
\label{prop:benchmark}
The present necessity / uniqueness theorem preserves the benchmark one-metric package \eqref{eq:fixedoutputs}. In particular, the operative metric remains exactly \(\CompletedMetric\), the ordinary matter route is unchanged, there is no second operative metric, and the low-energy matter law is not enlarged.
\end{proposition}

\begin{proof}
This is part of benchmark faithfulness in Definitions \ref{def:route}(7) and \ref{def:compatible}(5). The present note removes only the remaining body-level presentational freedom inside the current image-core frontier.
\end{proof}

\section{Conclusion}

The positive result of this note is narrower than a completed substrate derivation of GR but sharper than another routing reminder. On the current section-centered hidden-response route, the canonical image body is now proved to be the unique benchmark-facing body compatible with the fixed current displacement body, the fixed surjective projection, the fixed coercive response operator, the fixed exact-shell image, and the fixed current body-response route.

That matters because it spends the remaining body-level presentational freedom inside the current current-body-image-core frontier. Any benchmark-faithful compatible body is forced to be exactly the canonical minimal-lift image of the fixed current displacement body. The earlier projection-operator-core collapse theorem remains preserved main-line history because it first moved the live burden below the larger full-body presentation. The present necessity / uniqueness theorem then sharpens that smaller frontier by proving that even the image-body constituent itself is no longer independently variable at benchmark scope.

The scope remains honest. The one operative metric is still exactly \(\CompletedMetric\), the ordinary matter route is unchanged, there is no second operative metric, and the low-energy matter law is not enlarged. The next open burden is therefore no longer to restate or relocate the canonical image body itself. It is to derive or uniquely force the restricted projection / canonical-lift relation, the exact-shell image, the response operator restricted to the canonical image, or the full image core from still more primitive explicit substrate structure while preserving the same benchmark one-metric package and the same claim boundary.

\end{document}
